\chapter{Step 05: A Mini-ALU}
\label{ch:step05}

% Listings with line numbers, for listings whose lines the text refers to.
\lstdefinestyle{numbered}{numbers=left,numberstyle=\tiny\color{codecomment},numbersep=6pt}

% Inline code that contains a percent sign, such as a register name \ccpct{%1}.
\providecommand{\ccpct}{\begingroup\catcode`\%=12 \ccpctaux}
\providecommand{\ccpctaux}[1]{\cc{#1}\endgroup}

With step 05 the board starts to compute. The design of \cc{05_mini_alu} is
the Mini-ALU that we build in CircuitVerse in the course HPC0 (worksheets 4
and 6): four 8-bit registers, one of them the zero register, an
adder/subtractor, four status flags and a decoder. It executes 8-bit
instructions that we set on the DIP switch, one per press of a button. The
new pieces of SystemVerilog are the array, four registers under one name
selected by a signal, and the replication \cc{{8{sub}}}. And for the first
time we react to a button press rather than to the level of the button, with
the synchroniser and debouncer of Exercise~\ref{ex:step04-count}.

\section{The Mini-ALU from the Course}
\label{sec:step05-course}

Let us first recall what we built in CircuitVerse, so that we can compare it
with the text later.

\textbf{The instruction format.} An instruction is a bit pattern
$(\mathit{instr}_7, \dots, \mathit{instr}_0)$, split into four fields of two
bits each:
\begin{center}
\begin{tabular}{cccc}
\toprule
$\mathit{instr}_7\,\mathit{instr}_6$ & $\mathit{instr}_5\,\mathit{instr}_4$ &
$\mathit{instr}_3\,\mathit{instr}_2$ & $\mathit{instr}_1\,\mathit{instr}_0$ \\
\midrule
Op & X & Y & Z \\
\bottomrule
\end{tabular}
\end{center}
Op is the operation code, X, Y and Z name one of the four registers \ccpct{%0} to
\ccpct{%3}, and XY is the 4-bit pattern $(\mathit{instr}_5, \dots,
\mathit{instr}_2)$. With $u(\cdot)$ for the unsigned integer that a bit pattern
represents, worksheet 4 defines
\begin{center}
\begin{tabular}{cl}
\toprule
Op & operation \\
\midrule
00 & $u(\mathrm{XY}) \to u(\%Z)$ \\
01 & $(u(\mathrm{XY}) \cdot 2^4 + u(\%Z)) \bmod 2^8 \to u(\%Z)$ \\
10 & $(u(\%X) + u(\%Y)) \bmod 2^8 \to u(\%Z)$ \\
11 & $(u(\%X) - u(\%Y)) \bmod 2^8 \to u(\%Z)$ \\
\bottomrule
\end{tabular}
\end{center}
The first two operations together build any 8-bit value in a register: Op =
00 sets the low nibble and clears the high one, Op = 01 then adds the high
nibble.

\textbf{The parts.} The CircuitVerse Mini-ALU consists of
\begin{itemize}
\item a \emph{register file} with four 8-bit registers, one write port
  (\cc{data_in}, \cc{addr_in}) and two read ports (\cc{addr1_out},
  \cc{addr2_out}). A demultiplexer selects the register that receives the clock
  edge, a multiplexer per read port selects the register that is read. \ccpct{%0}
  is a \emph{zero register}: it has no flip-flops at all, reading it selects a
  constant 0, and a value written to it is lost;
\item an 8-bit \emph{adder/subtractor}: one adder whose second operand $b$ is
  passed through XOR gates, so that \cc{add/sub} $= 1$ inverts it and also
  feeds the carry-in (the XOR trick: $b \xor 0 = b$, $b \xor 1 = \neg b$);
\item the \emph{status flags}, computed from the operands $a$, $b$, the result
  $s$ and the carry-out, and stored in four flip-flops:
  ZF $= 1$ if $s = 0$; CF, the carry-out for an addition and the inverted
  carry-out for a subtraction; SF $= s_7$; and OF $= 1$ if $a$ and the added
  operand ($b$, or $\neg b$ for a subtraction) have the same sign and $s$ has
  the other one;
\item modules \cc{ZeroExtend} and \cc{LeftShift} for the immediate values of
  Op = 00 and Op = 01;
\item the \emph{ALU}, which connects them: the first operand of the adder is
  either the immediate value \cc{data_in} or the register \cc{data1_out},
  chosen by \cc{use_data1_out}; the second operand is always \cc{data2_out};
  the result goes back to the register file;
\item the \emph{decoder}, which splits the instruction with splitters and
  produces the control signals.
\end{itemize}
The only inputs left at the end are the instruction and the clock. The decoder
of worksheet 6 connects \cc{addr_in} to Z and \cc{addr1_out} to X, and
chooses the rest by Op:
\begin{center}
\begin{tabular}{cllcc}
\toprule
Op & \cc{data_in} & \cc{addr2_out} & \cc{add/sub} & \cc{use_data1_out} \\
\midrule
00 & \cc{ZeroExtend}(XY) & 00, so $b = u(\%0) = 0$ & 0 & 0 \\
01 & \cc{LeftShift}(XY) & Z & 0 & 0 \\
10 & (not used) & Y & 0 & 1 \\
11 & (not used) & Y & 1 & 1 \\
\bottomrule
\end{tabular}
\end{center}
So every instruction is an addition or a subtraction $a \pm b$, even the two
that load a value. (The table of operations in the decoder step of worksheet 6
writes Op = 11 as $(-u(\%X) + u(\%Y)) \bmod 2^8$. The circuit, with $a$ from
X and $b$ from Y, computes $u(\%X) - u(\%Y)$, as worksheet 4 and the final
table of worksheet 6 say, and so does our code.)

\section{The Description in SystemVerilog}
\label{sec:step05-sv}

The file \cc{test.sv} describes exactly this circuit
(Figure~\ref{fig:step05-datapath}), except for the clock, which we discuss in
Section~\ref{sec:step05-button}. We go through it in pieces.

\begin{figure}[ht]
\centering
\begin{tikzpicture}[x=1mm,y=1mm,font=\small,
    mux/.style={draw,rounded corners=2pt,fill=blue!7,minimum height=11mm,minimum width=5mm,inner sep=1pt,font=\scriptsize},
    op/.style={box,fill=black!5,minimum height=7mm}]
  % register file
  \node[box,fill=green!10,minimum width=17mm,minimum height=30mm,anchor=north west] (rf) at (0,30) {};
  \node[draw,dashed,minimum width=12mm,minimum height=5mm,font=\scriptsize] at (8.5,25) {\texttt{\%0} = 0};
  \foreach \i/\y in {1/18.5,2/12,3/5.5} {
    \node[draw,fill=white,minimum width=12mm,minimum height=5mm,font=\scriptsize] at (8.5,\y) {\texttt{\%\i}};
  }
  \node[note,rotate=90,anchor=south] at (rf.west) {\cc{register}};
  % read ports
  \node[mux] (mx) at (28,24) {X};
  \node[mux,minimum width=9mm] (my) at (28,6) {\cc{addr2}};
  \draw[wire,->] (rf.east |- mx) -- (mx.west);
  \draw[wire,->] (rf.east |- my) -- (my.west);
  % operand a
  \node[mux] (ma) at (46,22) {};
  \node[note,above=0.5mm of ma] {\cc{use_data1}};
  \node[note,anchor=east] (imm) at (42,19) {\cc{imm}};
  \draw[wire,->] (mx.east) -- (mx.east -| ma.west);
  \draw[wire,->] (imm.east) -- (imm.east -| ma.west);
  % operand b
  \node[op] (xor) at (46,6) {$\xor$ \cc{sub}};
  \draw[wire,->] (my.east) -- (xor.west);
  % adder
  \node[op,minimum height=26mm,minimum width=9mm] (add) at (64,15) {$+$};
  \node[note,below=0.5mm of add] {carry-in \cc{sub}};
  \draw[wire,->] (ma.east) -- (ma.east -| add.west);
  \draw[wire,->] (xor.east) -- (xor.east -| add.west);
  % flags
  \node[box,fill=green!10,minimum width=16mm] (fl) at (90,6) {ZF CF\\SF OF};
  \node[note,below=0.5mm of fl] {\cc{LED1}--\cc{LED4}};
  \draw[wire,->] (add.east |- fl) -- node[above,note]{\cc{s}, \cc{cout}} (fl.west);
  % result back to the register file
  \draw[wire,->] (add.east |- ma) -- (76,22) -- (76,38)
    -- node[above,note]{\cc{s}, written into \texttt{\%Z}} (8.5,38) -- (rf.north);
  % display
  \node[mux] (md) at (28,-12) {};
  \node[note,below=0mm of md] {\cc{BTN3}, \cc{BTN2}};
  \draw[wire,->] (rf.south) |- (md.west);
  \node[note,anchor=north east,align=right] at (7.5,-1) {clock \cc{CLK}\\enable: \cc{exec}, \cc{z} $\ne 0$};
  \node[box,fill=yellow!12] (hex) at (58,-12) {\cc{dev_hex}};
  \draw[wire,->] (md.east) -- (hex.west);
\end{tikzpicture}
\caption{The circuit described by \cc{05_mini_alu/test.sv}. The decoder (not
drawn) computes \cc{imm}, \cc{addr2}, \cc{sub} and \cc{use_data1} from the
instruction. Every read multiplexer gives 0 for address 0. At a rising edge of
\cc{CLK} in which \cc{exec} is 1, the sum \cc{s} is written into register Z
(unless Z is \texttt{\%0}) and the flags are stored.}
\label{fig:step05-datapath}
\end{figure}

\svpart{style=numbered,firstline=19,lastline=35,firstnumber=19}{05_mini_alu/test.sv}

\textbf{The ports.} \cc{A0} to \cc{A7} go to the 7-segment display, \cc{B0} to
\cc{B7} come from the DIP switch, as in Chapter~\ref{ch:step04}. \cc{BTN1}
executes the instruction, \cc{BTN2} and \cc{BTN3} select the register that the
display shows. New are the outputs \cc{LED1} to \cc{LED4} for the flags; they
are four of the five LEDs on PMOD 2, next to the buttons.

\textbf{Decoding the instruction.} Line 32 collects the switches into
\cc{instr}, with \cc{B0} as the most significant bit \cc{instr[7]}. So the
instruction reads in the order \cc{B0}, \dots, \cc{B7}, in the order in which
we write it down. Line 35 then splits it. A concatenation on the \emph{left}
side of an assignment (Chapter~\ref{ch:step03}) distributes the bits of the
right side, from the left: \cc{op} gets \cc{instr[7:6]}, \cc{x} gets
\cc{instr[5:4]}, \cc{y} gets \cc{instr[3:2]}, \cc{z} gets \cc{instr[1:0]}.
This is exactly the \cc{Instruction} module of worksheet 4: a splitter,
nothing but wires. The field XY needs no name of its own; where we need it, we
write \cc{{x, y}}.

\subsection*{The register file}

\svpart{style=numbered,firstline=59,lastline=61,firstnumber=59}{05_mini_alu/test.sv}

\textbf{An array of registers.} Line 61 declares an \term{array}: four
elements \cc{register[0]} to \cc{register[3]}, each a vector of 8 bits. The
range in front of the name, \cc{[7:0]}, gives the bits of one element; the
size after the name, \cc{[4]}, gives the number of elements and is short for
\cc{[0:3]}. (sv2v writes it like that: \cc{reg [7:0] register [0:3];}.)
\cc{register[2]} is an 8-bit vector, and \cc{register[2][3:0]} is its low
nibble. Like every \cc{logic}, the array says nothing yet about hardware. That
follows from its use, and it is used in two ways.

\textbf{Reading with a signal as index is a multiplexer.} In
\cc{register[x]} the index is not a constant but the 2-bit signal \cc{x}. The
expression stands for whichever of the four registers \cc{x} selects, that is,
a 4:1 multiplexer for 8-bit words with select input \cc{x}. The design has
three of these \term{read ports}: \cc{register[x]} for the first operand,
\cc{register[addr2]} for the second, and \cc{register[{BTN3, BTN2}]} for the
display. The register file of the course has two, plus the outputs
\ccpct{%0} to \ccpct{%3} for its four displays.

\textbf{The zero register.} Each read port is followed by an \cc{if} that
replaces the value by 0 when the address is 0, for example
\begin{lstlisting}
b = register[addr2];
if (addr2 == 2'b00) begin
    b = 8'h00;
end
\end{lstlisting}
(lines 103--106, the last assignment wins, Chapter~\ref{ch:step03}). This is
the multiplexer of the course with a constant 0 at its input 0. Since no read
port ever delivers \cc{register[0]}, and the write below skips it, nothing in
the design depends on \cc{register[0]}, and yosys removes its flip-flops: as
in the course, \ccpct{%0} has no storage.

\textbf{Writing with a signal as index selects the register.} In the
\cc{always_ff} block further down, \cc{register[z] <= s} says: at the edge, the
register with number \cc{z} takes the new value, and the other registers keep
theirs. In hardware, all flip-flops share one clock. Each register gets an
\term{enable} input: at the edge, a flip-flop with enable takes over its data
input only if the enable is 1, otherwise it keeps its value. A decoder computes
the enables from \cc{z}: register $i$ is enabled exactly when
$u(\mathtt{z}) = i$ (and, as we will see, \cc{exec} is 1). The
demultiplexer of worksheet 6 does the same job by passing the clock edge only
to the selected register. On an FPGA we do not switch clocks with logic, since
a clock should arrive at all flip-flops at the same moment; the enable
achieves the same result with a clean clock.

\subsection*{The decoder}

\svpart{style=numbered,firstline=63,lastline=86,firstnumber=63}{05_mini_alu/test.sv}

This is the table of Section~\ref{sec:step05-course}, written as a
\cc{case} in \cc{always_comb} (combinational logic, Chapter~\ref{ch:step03}).
\cc{imm} is the \cc{data_in} of the course: \cc{{4'h0, x, y}} is
\cc{ZeroExtend}, a concatenation with four zero bits on the left, and
\cc{{x, y, 4'h0}} is \cc{LeftShift}, four zero bits on the right, so
$u(\mathtt{\{x, y, 4'h0\}}) = u(\mathrm{XY}) \cdot 2^4$. Neither needs a
single gate. \cc{addr2} is the \cc{addr2_out} of the course, \cc{sub} is
\cc{add/sub}, and \cc{use_data1} is \cc{use_data1_out}. The \cc{default}
branch covers Op = 10 and 11, which differ only in \cc{sub}. Every branch
assigns both outputs, so no latch is needed (Chapter~\ref{ch:step03}).

\subsection*{The adder/subtractor}

\svpart{style=numbered,firstline=88,lastline=110,firstnumber=88}{05_mini_alu/test.sv}

Lines 92--107 are the two operand multiplexers of
Figure~\ref{fig:step05-datapath}, with the zero register built into the read
ports. Lines 109 and 110 are the adder/subtractor of worksheet 4.

\textbf{Replication.} \cc{{8{sub}}} is the bit \cc{sub} repeated eight
times, a \term{replication}: eight wires connected to the same signal. So
\cc{b ^ {8{sub}}} is a row of eight XOR gates, each with \cc{sub} at one
input, and \cc{b_xor} is $b$ for an addition and $\neg b$ for a subtraction.

\textbf{The carry-out.} The adder in line 110 has 9-bit operands: \cc{{1'b0,
a}} puts a zero bit in front of \cc{a}, and \cc{{8'h00, sub}} is the carry-in
as a 9-bit number. The 9-bit sum is split into its top bit \cc{cout} and the
8 bits \cc{s}. So
\[
  2^8 \cdot \mathit{cout} + u(s) = u(a) + u(b \xor \mathit{sub}\cdots\mathit{sub})
  + \mathit{sub},
\]
which for a subtraction is $u(a) + (2^8 - 1 - u(b)) + 1 = 2^8 + u(a) - u(b)$.
Thus $u(s) = (u(a) - u(b)) \bmod 2^8$, the two's-complement subtraction of the
course, and $\mathit{cout} = 1$ exactly when $u(a) \ge u(b)$. The \cc{+}
operator describes the whole carry chain; we do not write out the full adders.

\subsection*{The status flags}

\svpart{style=numbered,firstline=112,lastline=127,firstnumber=112}{05_mini_alu/test.sv}

The flip-flops \cc{ZF}, \cc{CF}, \cc{SF}, \cc{OF} are the
\cc{Status-Flags-FF} of worksheet 6, and the expressions are those of
worksheet 4:
\begin{itemize}
\item \cc{s == 8'h00} is a comparator with a constant, the NOR of all eight
  bits of $s$;
\item \cc{cout ^ sub} is CF. For an addition the carry-out says that the sum
  does not fit into 8 bits. For a subtraction we just saw that the carry-out
  is 1 exactly when there is \emph{no} borrow, $u(a) \ge u(b)$, so it has to be
  inverted, which is the XOR trick once more;
\item \cc{s[7]} is SF;
\item for OF, \cc{b_xor[7]} is the sign bit of the operand that is actually
  added, $b$ or $\neg b$. Signed overflow happens exactly when $a$ and this
  operand are both negative and $s$ is not, or both non-negative and $s$ is
  negative.
\end{itemize}
Line 127 shows the flags on the LEDs: \cc{LED1} is ZF, \cc{LED2} CF,
\cc{LED3} SF and \cc{LED4} OF.

Two details follow the course. First, the flags are stored for every executed
instruction, also for Op = 00 and Op = 01: the adder is used for them too.
After \cc{00001101}, which computes $3 + u(\%0) = 3 + 0$, all four flags are
0. Second, the flags do not depend on the destination. The \cc{if (z !=
2'b00)} only guards the write into the register file; the flags are assigned
outside of it. An instruction with destination \ccpct{%0} computes $s$ as
always, stores the flags, and discards $s$. This is the answer to the
question of worksheet 6, and it follows from the circuit, not from a special
rule.

\textbf{The display} shows \cc{register[{BTN3, BTN2}]}, with 0 for
\ccpct{%0}: with no button pressed \ccpct{%0}, with \cc{BTN2} \ccpct{%1},
with \cc{BTN3} \ccpct{%2}, with both \ccpct{%3}. The module \cc{dev_hex} is
the one of Chapter~\ref{ch:step03}.

\begin{exercise}
\label{ex:step05-decode}
Which instruction computes $(u(\%2) - u(\%1)) \bmod 2^8 \to u(\%3)$? Write
down the eight switch positions \cc{B0}, \dots, \cc{B7}. If \ccpct{%1} = 3 and
\ccpct{%2} = 2, what does the display show for \ccpct{%3} afterwards, and
which LEDs are on?
\end{exercise}

\begin{exercise}
\label{ex:step05-flags}
Write a program that puts \cc{80} (hexadecimal) into \ccpct{%1}, and then
computes $(u(\%1) + u(\%1)) \bmod 2^8 \to u(\%1)$. Predict the register and
the four flags after each instruction.
\end{exercise}

\section{The Button}
\label{sec:step05-button}

In the course, pressing the button produces the rising clock edge at which
the Mini-ALU executes the instruction. On the board, the 12\,MHz clock
\cc{CLK} is the only clock (Section~\ref{sec:step03-circuit}): all flip-flops
of the design, those of \cc{dev_hex}, of the register file and of the flags,
form one \term{clock domain}, the flip-flops driven by the same clock. The
button produces a signal \cc{exec} instead, which is 1 for exactly one clock
cycle per press, and the register file and the flags are enabled by it
(line 115). For the Mini-ALU, a press is then exactly one rising edge of the
course.

Why not simply use \cc{BTN1} as the clock? In the solution of
Exercise~\ref{ex:step02-button} we saw that a button is a bad clock: its
contacts bounce and produce several edges per press. For the Mini-ALU this
would mean that an instruction is sometimes executed several times.
Executing \cc{00001101} twice puts 3 into \ccpct{%1} twice, which does no
harm, but \cc{10011001}, $(u(\%1) + u(\%2)) \bmod 2^8 \to u(\%1)$, would add
\ccpct{%2} two or three times. The same happens if we only synchronise the
button and look for its rising edge, as Exercise~\ref{ex:step04-count}
showed: each bounce is a rising edge. We need a debouncer.

\svpart{style=numbered,firstline=37,lastline=57,firstnumber=37}{05_mini_alu/test.sv}

This is the circuit of the solution of Exercise~\ref{ex:step04-count}:
\begin{itemize}
\item \cc{btn_s1} and \cc{btn_s2} are the synchroniser
  (Section~\ref{sec:step04-button}). Only \cc{btn_s2} is used.
\item \cc{btn} is the debounced button, the value we have accepted.
  \cc{btn_cnt} is a 17-bit counter, which counts the cycles in which the
  synchronised button differs from \cc{btn}, and is reset to 0 as soon as both
  agree again. Only when they have differed for $2^{17}$ cycles in a row, that
  is, when \cc{btn_cnt} reaches \cc{17'h1ffff} and \cc{btn_s2} still differs,
  is the new value accepted. $2^{17}$ cycles at 12\,MHz are about 11\,ms,
  longer than the bouncing of a button, and short enough that we do not notice
  the delay. During the bouncing, the button keeps returning to its old value
  for a moment, which resets the counter; so the counter only reaches its end
  once the button is stable.
\item \cc{exec} is set in the same cycle in which \cc{btn} changes, to the new
  value. So it becomes 1 when a press is accepted, and not when a release is
  accepted. The default assignment \cc{exec <= 0} sets it back to 0 one cycle
  later.
\end{itemize}
As next-state functions, with $d = 1$ if \cc{btn_s2} $\ne$ \cc{btn} and $e = 1$
if additionally $u(\mathit{btn\_cnt}) = 2^{17} - 1$:
\[
  \mathit{btn\_cnt}' = d \cdot \bigl((u(\mathit{btn\_cnt}) + 1) \bmod 2^{17}\bigr), \qquad
  \mathit{btn}' = \mathit{btn} \xor e, \qquad
  \mathit{exec}' = e \wedge \mathit{btn\_s2} .
\]
Loading a value and keeping it, as in step 04, needed none of this, because
loading is idempotent (Section~\ref{sec:step04-button}). Executing an
instruction is not.

\begin{exercise}
\label{ex:step05-debounce}
How long after the button has stopped bouncing is the instruction executed?
What happens if you press the button only for 5\,ms? And if you press it
again 5\,ms after releasing it?
\end{exercise}

\section{What yosys Makes of It}
\label{sec:step05-yosys}

yosys first recognises the array as a \term{memory}, an array that is read and
written with a signal as index. It then decides how to build it:

\outputfile{step05_memory_map.txt}

``FF mapping'' means flip-flops, as opposed to one of the block RAMs of the
FPGA that we will meet in Chapter~\ref{ch:step07}. A block RAM would be far too
large for 32 bits, and above all it cannot do what we need: it has one read
port, and it delivers the data only after a clock edge. Our register file has
three read ports that deliver their values at once. So yosys builds four 8-bit
registers (``4 \$dff cells of width 8''), multiplexers for the read ports, and
the write logic. Later it notices that nothing depends on \cc{register[0]} and
removes it. After mapping to the building blocks of the iCE40:

\outputfile{step05_yosys_stat.txt}

(\cc{$scopeinfo} is bookkeeping about the module hierarchy, not hardware.)
There are 62 flip-flops. The 28 \cc{SB_DFFE}, with enable, are the registers
\ccpct{%1} to \ccpct{%3} (24) and the four flags, all enabled by \cc{exec}.
The 18 \cc{SB_DFFSR}, with synchronous reset, are \cc{btn_cnt}, reset when
the button agrees with \cc{btn}, and \cc{exec}. Of the 9 \cc{SB_DFF}, three
are \cc{btn_s1}, \cc{btn_s2} and \cc{btn}, and six belong to \cc{dev_hex},
as do the 7 \cc{SB_DFFESS} (Section~\ref{sec:step03-yosys}). Of the 26
\cc{SB_CARRY}, 3 belong to the counter of \cc{dev_hex}, 15 to \cc{btn_cnt},
and 8 to our adder/subtractor: one carry chain for all four operations, plus
one for the carry-out \cc{cout}.

nextpnr reports:

\outputfile{step05_pnr.txt}

One global buffer, since there is one clock. The slowest path runs from a
register through a read multiplexer and the operand logic into the carry chain
of the adder, and from its end back into a register: 21\,ns, so 47\,MHz, far
above our 12\,MHz. The switches are not synchronised; they are read only in
the cycle in which \cc{exec} is 1, and they do not change while we press the
button.

The DIP switch is on PMOD 1B, and \cc{icebreaker.pcf} names its eight pins
\cc{B0} to \cc{B7}; the LEDs and buttons of PMOD 2 are \cc{LED1} to \cc{LED5}
and \cc{BTN1} to \cc{BTN3}. So all 24 ports of \cc{test} have pins, and
nextpnr places them without warnings (Section~\ref{sec:step02-pins}).

\section{Simulation}
\label{sec:step05-sim}

The testbench \cc{tb_test.cpp} executes a program and checks every result.
Each instruction comes with the predicted registers and flags:

\cpppart{firstline=63,lastline=84}{05_mini_alu/tb_test.cpp}

The first four instructions are the program of worksheet 4, the fifth is from
worksheet 6, and the rest test the flags: the sixth is a subtraction with
destination \ccpct{%0}, the eighth an addition with a signed overflow,
$114 + 114 = 228$, which as a signed 8-bit number is $-28$.

For each instruction, the testbench sets the switches and presses and releases
\cc{BTN1} like a real button: the signal bounces four times within 2\,ms, and
then stays for $2^{17} + 100$ clock cycles. As in Chapter~\ref{ch:step04}, the
inputs change half a cycle before a rising edge. The testbench counts the
cycles in which \cc{exec} is 1, which must be exactly one per instruction.
Then it reads the four registers the way we do on the board: it sets
\cc{BTN2} and \cc{BTN3} and looks at the value \cc{hex_val} that \cc{dev_hex}
displays. To look at signals inside the design, the testbench includes the
header \cc{Vtest___024root.h}; then \cc{dut->rootp->test__DOT__exec} is the
signal \cc{exec} of the module \cc{test}, and
\cc{test__DOT__dev_hex0__DOT__hex_val} is \cc{hex_val} in the instance
\cc{dev_hex0}. The flags it reads from the outputs \cc{LED1} to \cc{LED4}.

\outputfile{step05_make_sim.txt}

Before an instruction writes them, \ccpct{%2} and \ccpct{%3} show
\cc{2F}: the simulation gives flip-flops without an initial value a random
value; on the board they start with 0. Those values are not checked. (The random values need the line
\cc{Verilated::commandArgs(argc, argv);} in the testbench, see
Appendix~\ref{app:tools}.) The simulation runs for about three million clock
cycles, a quarter of a second of the board, and takes less than a second. To
keep the waveform file small, the testbench writes only the first 20 cycles
of each press and each release, and the 20 cycles after each \cc{exec}, into
\cc{waveform.vcd}.

\begin{exercise}
\label{ex:step05-zero}
In the program above, \cc{11100100} has destination \ccpct{%0} and still sets
CF and SF. How would you change \cc{test.sv} so that an instruction with
destination \ccpct{%0} leaves the flags alone? Why is it useful that it does
not? Hint: think of comparing two registers.
\end{exercise}

\section{What Is Still Missing}

The Mini-ALU executes one instruction per button press, and we have to set
every instruction by hand. A computer fetches its instructions from a memory,
and to get a program into that memory, we need a way to talk to the board from
our computer. The next step does this: it connects the FPGA to the serial line
of the USB chip. The register file comes back in Chapter~\ref{ch:step09}, with
sixteen registers of 64 bits instead of four of eight, and again with a zero
register.

\solutions

\solution
We need Op = 11, X = 10, Y = 01 and Z = 11, so the instruction is
\cc{11100111}: \cc{B0} = 1, \cc{B1} = 1, \cc{B2} = 1, \cc{B3} = 0, \cc{B4} =
0, \cc{B5} = 1, \cc{B6} = 1, \cc{B7} = 1. With \ccpct{%1} = 3 and \ccpct{%2}
= 2 it computes $(2 - 3) \bmod 256 = 255$, so the display shows \cc{FF} for
\ccpct{%3}. ZF = 0. The adder computes $2 + 252 + 1 = 255$ without carry-out,
so there is a borrow and CF = 1. SF = 1. OF = 0: $a = 2$ is non-negative and
the added operand $\neg b$ = \cc{11111100} is negative, and the sum of two
numbers with different signs never overflows. So \cc{LED2} and \cc{LED3} are
on.

\solution
\cc{00000001} loads 0 into \ccpct{%1} (ZF = 1, since the sum is 0, and the
other flags 0), \cc{01100001} adds
$u(1000) \cdot 2^4 = 128$, so \ccpct{%1} = \cc{80} (ZF = 0, CF = 0, SF = 1, OF =
0: $0 + (-128)$ does not overflow). Then \cc{10010101} computes
$128 + 128 = 256$, so \ccpct{%1} = \cc{00}, with ZF = 1, CF = 1 (unsigned
overflow), SF = 0 and OF = 1: two negative operands, $-128 + (-128) = -256$,
give a non-negative result. (If \ccpct{%1} was 0 already, the first
instruction is not needed.) Putting the three instructions with these
predictions into the table of the testbench confirms them.

\solution
\cc{btn} changes when the synchronised button has differed from it for
$2^{17}$ cycles in a row, so the instruction is executed about $2^{17}$
cycles, 11\,ms, after the last bounce. (A few cycles more: two for the
synchroniser, and \cc{exec} is a register itself. In the simulation, \cc{exec}
is 1 exactly $2^{17} + 2$ cycles after the last change of \cc{BTN1}.) A press of 5\,ms is ignored: the counter is reset when the button is
released, before it reaches its end. A press 5\,ms after a release is also
ignored: the release has not been accepted yet, \cc{btn} is still 1, and the
new press only means that the button agrees with \cc{btn} again. On the
board, both are far shorter than a normal press.

\solution
Move the four flag assignments into the \cc{if (z != 2'b00)} block. But an
instruction that computes flags and discards its result is useful:
\cc{11011000}, for example, computes $u(\%1) - u(\%2)$ and only keeps the
flags. Afterwards ZF says whether the registers are equal and CF whether
$u(\%1) < u(\%2)$, without destroying a register. Many processors have such a
\emph{compare} instruction, and with a zero register we get it for free, as
worksheet 6 points out.
